% ============================================================================
%  SECCIÓN 4.1: IDENTIFICACIÓN DE LA LÓGICA DE NEGOCIO
% ============================================================================

\section{Identificaci\'on de la l\'ogica de negocio}

En el PDF 3 se construyeron las principales interfaces de la aplicaci\'on y se
estableci\'o la navegaci\'on entre ellas. En este avance, la aplicaci\'on comienza
a responder a las acciones del usuario. El cambio fundamental es que la
aplicaci\'on deja de ser una colecci\'on est\'atica de pantallas y comienza a tomar
decisiones y responder a los datos introducidos por el usuario.

\subsection{Pantallas que requieren l\'ogica}

No todas las pantallas de la aplicaci\'on requieren el mismo nivel de
interacci\'on. A continuaci\'on se identifican las pantallas que necesitan
l\'ogica de negocio y el tipo de interacci\'on que manejan:

\begin{table}[H]
  \centering
  \caption{Pantallas con l\'ogica de negocio}
  \contenidoConFuente{%
    \begin{tablaAPA}
      \begin{tabular}{@{}p{3.8cm}p{11.8cm}@{}}
        \toprule
        \textbf{Pantalla} & \textbf{Interacci\'on} \\
        \midrule
        Inicio de sesi\'on & Formulario con validaci\'on de credenciales,
        manejo de errores del servidor y indicador de carga. \\
        \addlinespace
        Registro & Formulario con validaci\'on de cuatro campos,
        confirmaci\'on de contrase\~na y mensaje de \'exito. \\
        \addlinespace
        Cat\'alogo (Buscar) & B\'usqueda por nombre con debounce,
        filtros por tipo y deporte, limpieza de filtros. \\
        \addlinespace
        Detalle del escenario & Toggle de favorito, subida de imagen
        (solo admin/gestor), visualizaci\'on de eventos. \\
        \addlinespace
        Favoritos & Lista din\'amica de escenarios guardados,
        pull-to-refresh. \\
        \addlinespace
        Perfil & Cierre de sesi\'on con di\'alogo de confirmaci\'on. \\
        \bottomrule
      \end{tabular}
    \end{tablaAPA}%
  }{Elaboraci\'on propia en base al an\'alisis del c\'odigo fuente.}
\end{table}

\subsection{Flujo de datos de la aplicaci\'on}

El flujo general de datos en la aplicaci\'on sigue el patr\'on propuesto por la
gu\'ia del avance:

\begin{center}
\begin{tikzpicture}[
  node distance=0.7cm,
  nodo/.style={rectangle, draw=black!60, rounded corners, minimum width=2.8cm,
    minimum height=0.7cm, align=center, font=\small},
  flecha/.style={-Stealth, thick}
]
  \node[nodo] (usuario) {Usuario};
  \node[nodo, below=of usuario] (interactua) {Interact\'ua};
  \node[nodo, below=of interactua] (app) {Aplicaci\'on};
  \node[nodo, below=of app] (procesa) {Procesa};
  \node[nodo, below=of procesa] (estado) {Actualiza estado};
  \node[nodo, below=of estado] (muestra) {Muestra resultado};
  \draw[flecha] (usuario) -- (interactua);
  \draw[flecha] (interactua) -- (app);
  \draw[flecha] (app) -- (procesa);
  \draw[flecha] (procesa) -- (estado);
  \draw[flecha] (estado) -- (muestra);
\end{tikzpicture}
\end{center}

En la pr\'actica, este flujo se implementa de la siguiente manera: el usuario
presiona un bot\'on o ingresa datos en un campo de texto; la aplicaci\'on captura
la acci\'on mediante un controlador o un evento \textit{onPress}; la l\'ogica de
negocio valida los datos, ejecuta la operaci\'on correspondiente (consulta a
Supabase, actualizaci\'on de estado local o navegaci\'on) y, finalmente, la interfaz
se actualiza para reflejar el nuevo estado de la aplicaci\'on.

\subsection{Operaciones de negocio identificadas}

A partir del an\'alisis de las pantallas, se identificaron las siguientes
operaciones de negocio principales:

\begin{itemize}
  \item \textbf{Autenticaci\'on:} inicio de sesi\'on, registro de nuevos usuarios
    y cierre de sesi\'on, gestionados a trav\'es del contexto de autenticaci\'on
    (\textit{AuthContext}).
  \item \textbf{Consulta de escenarios:} obtenci\'on de la lista de escenarios
    deportivos activos desde Supabase, con soporte para b\'usqueda y filtros.
  \item \textbf{Gesti\'on de favoritos:} adicion y eliminaci\'on de escenarios a la
    lista de favoritos del usuario, con sincronizaci\'on en tiempo real.
  \item \textbf{Detalle del escenario:} visualizaci\'on de informaci\'on completa,
    incluyendo imagen, deportes disponibles y eventos pr\'oximos.
  \item \textbf{Subida de imagen:} carga de fotograf\'ias de escenarios,
    restringida a usuarios con rol de administrador o gestor.
\end{itemize}
